\chapter{Analizador sintáctico}
\label{cap:analizador-sintactico}

El analizador sintáctico verifica que la secuencia de tokens respete la
estructura del lenguaje TinyS. La implementación se basa en un analizador
sintáctico predictivo descendente recursivo, que utiliza los conjuntos
de primeros y siguientes y un token de anticipación para seleccionar las 
producciones y comprobar los símbolos esperados. En esta etapa no se realizan 
controles de tipos ni se construye el árbol sintáctico abstracto.

% BORRADOR INICIAL: completar todas las transformaciones de la gramática,
% la gramática definitiva completa y los conjuntos FIRST y FOLLOW.
% Los ejemplos siguientes muestran transformaciones concretas, pero no
% constituyen todavía la documentación completa exigida para esta etapa.

\lstdefinestyle{bnftinys}{
  basicstyle=\ttfamily\small,
  columns=fullflexible,
  keepspaces=true,
  breaklines=true,
  breakatwhitespace=true,
  breakindent=2em,
  frame=single,
  rulecolor=\color{gray},
  showstringspaces=false,
  aboveskip=0.8em,
  belowskip=0.8em,
  literate=
    {á}{{\'a}}1 {é}{{\'e}}1 {í}{{\'i}}1 {ó}{{\'o}}1 {ú}{{\'u}}1
    {Á}{{\'A}}1 {É}{{\'E}}1 {Í}{{\'I}}1 {Ó}{{\'O}}1 {Ú}{{\'U}}1
    {ñ}{{\~n}}1 {Ñ}{{\~N}}1
    {⟨}{{$\langle$}}1 {⟩}{{$\rangle$}}1 {λ}{{$\lambda$}}1
}

\section{Gramática BNF Extendida}
\lstinputlisting[style=bnftinys]{gramaticas/01_ebnf_original.bnf}

\section{Transformación a BNF}
En esta etapa, la estrella de Kleene se reemplaza por nuevas reglas con una K mayúscula al final del nombre del no terminal. Por ejemplo, la regla
\begin{lstlisting}[style=bnftinys]
<Bloque> ::= { <Sentencia>* }
\end{lstlisting}
se transforma en:
\begin{lstlisting}[style=bnftinys]
<Bloque> ::= { <SentenciaK> }
<SentenciaK> ::= <Sentencia> <SentenciaK>
	| λ
\end{lstlisting}
Y para el caso de la opcionalidad "\texttt{?}", se crean nuevas reglas para reflejar ese comportamiento, prácticamente se duplica la regla, con la única diferencia de que
la que tenía la opcionalidad en una aparece y en otra no.
\lstinputlisting[style=bnftinys]{gramaticas/02_bnf.bnf}

\section{Eliminación de ambigüedad y redundancia}
\label{sec:eliminacion-ambiguedad}
En esta sección se muestran las transformaciones aplicadas a la gramática para eliminar ambigüedad y redundancia,
pero por la naturaleza de las sentencias if y else, no es posible eliminar completamente la ambigüedad.
\lstinputlisting[style=bnftinys]{gramaticas/03_sin_ambiguedad_y_redundancia.bnf}

\section{Eliminación de recursividad por izquierda}
En esta etapa, para eliminar la recursividad por izquierda, se crean nuevas reglas con una R mayúscula al final del nombre del no terminal. Por ejemplo, la regla
\begin{lstlisting}[style=bnftinys]
  <ExpOr> ::= <ExpOr> || <ExpAnd>
    | <ExpAnd>
\end{lstlisting}
se transforma en:
\begin{lstlisting}[style=bnftinys]<ExpOr> ::= <ExpAnd> <ExpOrR>
<ExpOr> ::= <ExpAnd> <ExpOrR>
<ExpOrR> ::= || <ExpAnd> <ExpOrR>
  | λ
\end{lstlisting}
\lstinputlisting[style=bnftinys]{gramaticas/04_sin_recursividad_por_izquierda.bnf}

\section{Factorización por izquierda}
En esta etapa, para realizar la factorización por izquierda, se crean nuevas reglas con una F mayúscula al final del nombre del no terminal. Por ejemplo, la regla
\begin{lstlisting}[style=bnftinys]
  <class> ::= class idclass { <AtributoK> }
    | class idclass <Herencia> { <AtributoK> }
\end{lstlisting}
se transforma en:
\begin{lstlisting}[style=bnftinys]
  <class> ::= class idclass <classF>

  <classF> ::= { <AtributoK> }
    |<Herencia> { <AtributoK> }
\end{lstlisting}
Además, se constató que ⟨Visibilidad⟩ y ⟨Forma-Método⟩ son reglas de renombramiento, por lo que se reemplazaron por sus producciones. La gramática final es la siguiente:
\lstinputlisting[style=bnftinys]{gramaticas/05_factorizacion_por_izquierda.bnf}


\section{Primeros}
\begingroup
\small
\renewcommand{\arraystretch}{1.2}
\setlength{\tabcolsep}{5pt}
\setlength{\LTleft}{0pt}
\setlength{\LTright}{0pt plus 1fill}
\begin{longtable}{|
>{\raggedright\arraybackslash}p{0.43\textwidth}|
>{\raggedright\arraybackslash}p{\dimexpr0.57\textwidth-4\tabcolsep-3\arrayrulewidth\relax}|}
\caption{Conjuntos de primeros}
\label{tab:primeros}\\
\hline
\textbf{No terminal} & \textbf{Primeros} \\ \hline
\endfirsthead
\hline
\textbf{No terminal} & \textbf{Primeros} \\ \hline
\endhead
\texttt{program} & \(\{\) \texttt{class}, \texttt{impl}, \texttt{start} \(\}\) \\ \hline
\texttt{Start} & \(\{\) \texttt{start} \(\}\) \\ \hline
\texttt{Lista-Definiciones} & \(\{\) \(\lambda\), \texttt{class}, \texttt{impl} \(\}\) \\ \hline
\texttt{class} & \(\{\) \texttt{class} \(\}\) \\ \hline
\texttt{classF} & \(\{\) \texttt{:}, \texttt{'\{'} \(\}\) \\ \hline
\texttt{AtributoK} & \(\{\) \(\lambda\), \texttt{Str}, \texttt{Bool}, \texttt{Int}, \texttt{idclass}, \texttt{Array}, \texttt{pub} \(\}\) \\ \hline
\texttt{Impl} & \(\{\) \texttt{impl} \(\}\) \\ \hline
\texttt{Herencia} & \(\{\) \texttt{:} \(\}\) \\ \hline
\texttt{MiembroK} & \(\{\) \texttt{.}, \texttt{st}, \texttt{fn}, \(\lambda\) \(\}\) \\ \hline
\texttt{Miembro} & \(\{\) \texttt{.}, \texttt{st}, \texttt{fn} \(\}\) \\ \hline
\texttt{Constructor} & \(\{\) \texttt{.} \(\}\) \\ \hline
\texttt{Atributo} & \(\{\) \texttt{Str}, \texttt{Bool}, \texttt{Int}, \texttt{idclass}, \texttt{Array}, \texttt{pub} \(\}\) \\ \hline
\texttt{Método} & \(\{\) \texttt{st}, \texttt{fn} \(\}\) \\ \hline
\texttt{MétodoF} & \(\{\) \texttt{idMetAt}, \texttt{void}, \texttt{Str}, \texttt{Bool}, \texttt{Int}, \texttt{idclass}, \texttt{Array} \(\}\) \\ \hline
\texttt{Bloque-Método} & \(\{\) \texttt{'\{'} \(\}\) \\ \hline
\texttt{Decl-Var-LocalesK} & \(\{\) \(\lambda\), \texttt{Str}, \texttt{Bool}, \texttt{Int}, \texttt{idclass}, \texttt{Array} \(\}\) \\ \hline
\texttt{Decl-Var-Locales} & \(\{\) \texttt{Str}, \texttt{Bool}, \texttt{Int}, \texttt{idclass}, \texttt{Array} \(\}\) \\ \hline
\texttt{Lista-Declaración-Variables} & \(\{\) \texttt{idMetAt} \(\}\) \\ \hline
\texttt{Lista-Declaración-VariablesF} & \(\{\) \(\lambda\), \texttt{','} \(\}\) \\ \hline
\texttt{Argumentos-Formales} & \(\{\) \texttt{(} \(\}\) \\ \hline
\texttt{Argumentos-FormalesF} & \(\{\) \texttt{)}, \texttt{Str}, \texttt{Bool}, \texttt{Int}, \texttt{idclass}, \texttt{Array} \(\}\) \\ \hline
\texttt{Lista-Argumentos-Formales} & \(\{\) \texttt{Str}, \texttt{Bool}, \texttt{Int}, \texttt{idclass}, \texttt{Array} \(\}\) \\ \hline
\texttt{Lista-Argumentos-FormalesF} & \(\{\) \(\lambda\), \texttt{','} \(\}\) \\ \hline
\texttt{Argumento-Formal} & \(\{\) \texttt{Str}, \texttt{Bool}, \texttt{Int}, \texttt{idclass}, \texttt{Array} \(\}\) \\ \hline
\texttt{Tipo-Método} & \(\{\) \texttt{void}, \texttt{Str}, \texttt{Bool}, \texttt{Int}, \texttt{idclass}, \texttt{Array} \(\}\) \\ \hline
\texttt{Tipo} & \(\{\) \texttt{Str}, \texttt{Bool}, \texttt{Int}, \texttt{idclass}, \texttt{Array} \(\}\) \\ \hline
\texttt{Tipo-Primitivo} & \(\{\) \texttt{Str}, \texttt{Bool}, \texttt{Int} \(\}\) \\ \hline
\texttt{Tipo-Referencia} & \(\{\) \texttt{idclass} \(\}\) \\ \hline
\texttt{Tipo-Arreglo} & \(\{\) \texttt{Array} \(\}\) \\ \hline
\texttt{SentenciaK} & \(\{\) \texttt{;}, \texttt{while}, \texttt{for}, \texttt{if}, \texttt{ret}, \texttt{(}, \texttt{'\{'}, \texttt{id}, \texttt{self}, \(\lambda\) \(\}\) \\ \hline
\texttt{Sentencia} & \(\{\) \texttt{;}, \texttt{while}, \texttt{for}, \texttt{if}, \texttt{ret}, \texttt{(}, \texttt{'\{'}, \texttt{id}, \texttt{self} \(\}\) \\ \hline
\texttt{SentenciaF1} & \(\{\) \texttt{else}, \(\lambda\) \(\}\) \\ \hline
\texttt{SentenciaF2} & \(\{\) \texttt{;}, \texttt{+}, \texttt{-}, \texttt{!}, \texttt{++}, \texttt{-{}-}, \texttt{nil}, \texttt{true}, \texttt{false}, \texttt{intLiteral}, \texttt{StrLiteral}, \texttt{(}, \texttt{self}, \texttt{idclass}, \texttt{new}, \texttt{id} \(\}\) \\ \hline
\texttt{Bloque} & \(\{\) \texttt{'\{'} \(\}\) \\ \hline
\texttt{Asignación} & \(\{\) \texttt{id}, \texttt{self} \(\}\) \\ \hline
\texttt{AccesoVar-Simple} & \(\{\) \texttt{id} \(\}\) \\ \hline
\texttt{AccesoVar-SimpleF} & \(\{\) \(\lambda\), \texttt{.}, \texttt{[} \(\}\) \\ \hline
\texttt{AccesoSelf-Simple} & \(\{\) \texttt{self} \(\}\) \\ \hline
\texttt{Encadenado-SimpleK} & \(\{\) \(\lambda\), \texttt{.} \(\}\) \\ \hline
\texttt{Encadenado-Simple} & \(\{\) \texttt{.} \(\}\) \\ \hline
\texttt{Sentencia-Simple} & \(\{\) \texttt{(} \(\}\) \\ \hline
\texttt{Expresión} & \(\{\) \texttt{+}, \texttt{-}, \texttt{!}, \texttt{++}, \texttt{-{}-}, \texttt{nil}, \texttt{true}, \texttt{false}, \texttt{intLiteral}, \texttt{StrLiteral}, \texttt{(}, \texttt{self}, \texttt{idclass}, \texttt{new}, \texttt{id} \(\}\) \\ \hline
\texttt{ExpOr} & \(\{\) \texttt{+}, \texttt{-}, \texttt{!}, \texttt{++}, \texttt{-{}-}, \texttt{nil}, \texttt{true}, \texttt{false}, \texttt{intLiteral}, \texttt{StrLiteral}, \texttt{(}, \texttt{self}, \texttt{idclass}, \texttt{new}, \texttt{id} \(\}\) \\ \hline
\texttt{ExpOrR} & \(\{\) \texttt{||}, \(\lambda\) \(\}\) \\ \hline
\texttt{ExpAnd} & \(\{\) \texttt{+}, \texttt{-}, \texttt{!}, \texttt{++}, \texttt{-{}-}, \texttt{nil}, \texttt{true}, \texttt{false}, \texttt{intLiteral}, \texttt{StrLiteral}, \texttt{(}, \texttt{self}, \texttt{idclass}, \texttt{new}, \texttt{id} \(\}\) \\ \hline
\texttt{ExpAndR} & \(\{\) \texttt{\&\&}, \(\lambda\) \(\}\) \\ \hline
\texttt{ExpIgual} & \(\{\) \texttt{+}, \texttt{-}, \texttt{!}, \texttt{++}, \texttt{-{}-}, \texttt{nil}, \texttt{true}, \texttt{false}, \texttt{intLiteral}, \texttt{StrLiteral}, \texttt{(}, \texttt{self}, \texttt{idclass}, \texttt{new}, \texttt{id} \(\}\) \\ \hline
\texttt{ExpIgualR} & \(\{\) \texttt{==}, \texttt{!=}, \(\lambda\) \(\}\) \\ \hline
\texttt{ExpCompuesta} & \(\{\) \texttt{+}, \texttt{-}, \texttt{!}, \texttt{++}, \texttt{-{}-}, \texttt{nil}, \texttt{true}, \texttt{false}, \texttt{intLiteral}, \texttt{StrLiteral}, \texttt{(}, \texttt{self}, \texttt{idclass}, \texttt{new}, \texttt{id} \(\}\) \\ \hline
\texttt{ExpCompuestaF} & \(\{\) \texttt{<}, \texttt{>}, \texttt{<=}, \texttt{>=}, \(\lambda\) \(\}\) \\ \hline
\texttt{ExpAd} & \(\{\) \texttt{+}, \texttt{-}, \texttt{!}, \texttt{++}, \texttt{-{}-}, \texttt{nil}, \texttt{true}, \texttt{false}, \texttt{intLiteral}, \texttt{StrLiteral}, \texttt{(}, \texttt{self}, \texttt{idclass}, \texttt{new}, \texttt{id} \(\}\) \\ \hline
\texttt{ExpAdR} & \(\{\) \texttt{+}, \texttt{-}, \(\lambda\) \(\}\) \\ \hline
\texttt{ExpMul} & \(\{\) \texttt{+}, \texttt{-}, \texttt{!}, \texttt{++}, \texttt{-{}-}, \texttt{nil}, \texttt{true}, \texttt{false}, \texttt{intLiteral}, \texttt{StrLiteral}, \texttt{(}, \texttt{self}, \texttt{idclass}, \texttt{new}, \texttt{id} \(\}\) \\ \hline
\texttt{ExpMulR} & \(\{\) \texttt{*}, \texttt{/}, \(\lambda\) \(\}\) \\ \hline
\texttt{ExpUn} & \(\{\) \texttt{+}, \texttt{-}, \texttt{!}, \texttt{++}, \texttt{-{}-}, \texttt{nil}, \texttt{true}, \texttt{false}, \texttt{intLiteral}, \texttt{StrLiteral}, \texttt{(}, \texttt{self}, \texttt{idclass}, \texttt{new}, \texttt{id} \(\}\) \\ \hline
\texttt{OpIgual} & \(\{\) \texttt{==}, \texttt{!=} \(\}\) \\ \hline
\texttt{OpCompuesto} & \(\{\) \texttt{<}, \texttt{>}, \texttt{<=}, \texttt{>=} \(\}\) \\ \hline
\texttt{OpAd} & \(\{\) \texttt{+}, \texttt{-} \(\}\) \\ \hline
\texttt{OpUnario} & \(\{\) \texttt{+}, \texttt{-}, \texttt{!}, \texttt{++}, \texttt{-{}-} \(\}\) \\ \hline
\texttt{OpMul} & \(\{\) \texttt{*}, \texttt{/} \(\}\) \\ \hline
\texttt{Operando} & \(\{\) \texttt{nil}, \texttt{true}, \texttt{false}, \texttt{intLiteral}, \texttt{StrLiteral}, \texttt{(}, \texttt{self}, \texttt{idclass}, \texttt{new}, \texttt{id} \(\}\) \\ \hline
\texttt{Literal} & \(\{\) \texttt{nil}, \texttt{true}, \texttt{false}, \texttt{intLiteral}, \texttt{StrLiteral} \(\}\) \\ \hline
\texttt{Primario} & \(\{\) \texttt{(}, \texttt{self}, \texttt{idclass}, \texttt{new}, \texttt{id} \(\}\) \\ \hline
\texttt{VarOMetodo} & \(\{\) \(\lambda\), \texttt{[}, \texttt{.}, \texttt{(} \(\}\) \\ \hline
\texttt{ExpresionParentizada} & \(\{\) \texttt{(} \(\}\) \\ \hline
\texttt{ExpresionParentizadaF} & \(\{\) \(\lambda\), \texttt{.} \(\}\) \\ \hline
\texttt{AccesoSelf} & \(\{\) \texttt{self} \(\}\) \\ \hline
\texttt{AccesoSelfF} & \(\{\) \(\lambda\), \texttt{.} \(\}\) \\ \hline
\texttt{AccesoVar} & \(\{\) \(\lambda\), \texttt{[}, \texttt{.} \(\}\) \\ \hline
\texttt{AccesoVarF} & \(\{\) \(\lambda\), \texttt{.} \(\}\) \\ \hline
\texttt{Llamada-Método} & \(\{\) \texttt{(} \(\}\) \\ \hline
\texttt{Llamada-MétodoF} & \(\{\) \(\lambda\), \texttt{.} \(\}\) \\ \hline
\texttt{Llamada-Método-Estático} & \(\{\) \texttt{idclass} \(\}\) \\ \hline
\texttt{Llamada-Conclassor} & \(\{\) \texttt{new} \(\}\) \\ \hline
\texttt{Llamada-ConclassorF1} & \(\{\) \texttt{idclass}, \texttt{Str}, \texttt{Bool}, \texttt{Int} \(\}\) \\ \hline
\texttt{Llamada-ConclassorF2} & \(\{\) \(\lambda\), \texttt{.} \(\}\) \\ \hline
\texttt{Argumentos-Actuales} & \(\{\) \texttt{(} \(\}\) \\ \hline
\texttt{Argumentos-ActualesF} & \(\{\) \texttt{)}, \texttt{+}, \texttt{-}, \texttt{!}, \texttt{++}, \texttt{-{}-}, \texttt{nil}, \texttt{true}, \texttt{false}, \texttt{intLiteral}, \texttt{StrLiteral}, \texttt{(}, \texttt{self}, \texttt{idclass}, \texttt{new}, \texttt{id} \(\}\) \\ \hline
\texttt{Lista-Expresiones} & \(\{\) \texttt{+}, \texttt{-}, \texttt{!}, \texttt{++}, \texttt{-{}-}, \texttt{nil}, \texttt{true}, \texttt{false}, \texttt{intLiteral}, \texttt{StrLiteral}, \texttt{(}, \texttt{self}, \texttt{idclass}, \texttt{new}, \texttt{id} \(\}\) \\ \hline
\texttt{Lista-ExpresionesF} & \(\{\) \texttt{','}, \(\lambda\) \(\}\) \\ \hline
\texttt{Encadenado} & \(\{\) \texttt{.} \(\}\) \\ \hline
\texttt{EncadenadoF} & \(\{\) \texttt{(}, \(\lambda\), \texttt{[}, \texttt{.} \(\}\) \\ \hline
\end{longtable}
\endgroup


\clearpage
\section{Siguientes}
\begingroup
\small
\renewcommand{\arraystretch}{1.2}
\setlength{\tabcolsep}{5pt}
\setlength{\LTleft}{0pt}
\setlength{\LTright}{0pt plus 1fill}
\begin{longtable}{|
>{\raggedright\arraybackslash}p{0.43\textwidth}|
>{\raggedright\arraybackslash}p{\dimexpr0.57\textwidth-4\tabcolsep-3\arrayrulewidth\relax}|}
\caption{Conjuntos de siguientes}
\label{tab:siguientes}\\
\hline
\textbf{No terminal} & \textbf{Siguientes} \\ \hline
\endfirsthead
\hline
\textbf{No terminal} & \textbf{Siguientes} \\ \hline
\endhead
\texttt{program} & \(\{\) \texttt{éxito} \(\}\) \\ \hline
\texttt{Lista-Definiciones} & \(\{\) \texttt{start} \(\}\) \\ \hline
\texttt{Start} & \(\{\) \texttt{éxito} \(\}\) \\ \hline
\texttt{Bloque-Método} & \(\{\) \texttt{éxito}, \texttt{.}, \texttt{st}, \texttt{fn}, \texttt{'\}'} \(\}\) \\ \hline
\texttt{class} & \(\{\) \texttt{class}, \texttt{impl}, \texttt{start} \(\}\) \\ \hline
\texttt{Impl} & \(\{\) \texttt{class}, \texttt{impl}, \texttt{start} \(\}\) \\ \hline
\texttt{classF} & \(\{\) \texttt{class}, \texttt{impl}, \texttt{start} \(\}\) \\ \hline
\texttt{AtributoK} & \(\{\) \texttt{'\}'} \(\}\) \\ \hline
\texttt{Herencia} & \(\{\) \texttt{'\{'} \(\}\) \\ \hline
\texttt{Atributo} & \(\{\) \texttt{Str}, \texttt{Bool}, \texttt{Int}, \texttt{idclass}, \texttt{Array}, \texttt{pub}, \texttt{'\}'} \(\}\) \\ \hline
\texttt{MiembroK} & \(\{\) \texttt{'\}'} \(\}\) \\ \hline
\texttt{Tipo} & \(\{\) \texttt{'\{'}, \texttt{idMetAt} \(\}\) \\ \hline
\texttt{Miembro} & \(\{\) \texttt{.}, \texttt{st}, \texttt{fn}, \texttt{'\}'} \(\}\) \\ \hline
\texttt{Método} & \(\{\) \texttt{.}, \texttt{st}, \texttt{fn}, \texttt{'\}'} \(\}\) \\ \hline
\texttt{Constructor} & \(\{\) \texttt{.}, \texttt{st}, \texttt{fn}, \texttt{'\}'} \(\}\) \\ \hline
\texttt{Argumentos-Formales} & \(\{\) \texttt{'\{'} \(\}\) \\ \hline
\texttt{Lista-Declaración-Variables} & \(\{\) \texttt{;} \(\}\) \\ \hline
\texttt{Visibilidad} & \(\{\) \texttt{Str}, \texttt{Bool}, \texttt{Int}, \texttt{idclass}, \texttt{Array} \(\}\) \\ \hline
\texttt{MétodoF} & \(\{\) \texttt{.}, \texttt{st}, \texttt{fn}, \texttt{'\}'} \(\}\) \\ \hline
\texttt{Forma-Método} & \(\{\) \texttt{fn} \(\}\) \\ \hline
\texttt{Tipo-Método} & \(\{\) \texttt{idMetAt} \(\}\) \\ \hline
\texttt{Decl-Var-LocalesK} & \(\{\) \texttt{;}, \texttt{while}, \texttt{for}, \texttt{if}, \texttt{ret}, \texttt{(}, \texttt{'\{'}, \texttt{id}, \texttt{self}, \texttt{'\}'} \(\}\) \\ \hline
\texttt{SentenciaK} & \(\{\) \texttt{'\}'} \(\}\) \\ \hline
\texttt{Decl-Var-Locales} & \(\{\) \texttt{Str}, \texttt{Bool}, \texttt{Int}, \texttt{idclass}, \texttt{Array}, \texttt{;}, \texttt{while}, \texttt{for}, \texttt{if}, \texttt{ret}, \texttt{(}, \texttt{'\{'}, \texttt{id}, \texttt{self}, \texttt{'\}'} \(\}\) \\ \hline
\texttt{Lista-Declaración-VariablesF} & \(\{\) \texttt{;} \(\}\) \\ \hline
\texttt{Argumentos-FormalesF} & \(\{\) \texttt{'\{'} \(\}\) \\ \hline
\texttt{Lista-Argumentos-Formales} & \(\{\) \texttt{)} \(\}\) \\ \hline
\texttt{Argumento-Formal} & \(\{\) \texttt{','}, \texttt{)} \(\}\) \\ \hline
\texttt{Lista-Argumentos-FormalesF} & \(\{\) \texttt{)} \(\}\) \\ \hline
\texttt{Tipo-Primitivo} & \(\{\) \texttt{idMetAt}, \texttt{[}, \texttt{'\{'} \(\}\) \\ \hline
\texttt{Tipo-Referencia} & \(\{\) \texttt{'\{'}, \texttt{idMetAt} \(\}\) \\ \hline
\texttt{Tipo-Arreglo} & \(\{\) \texttt{'\{'}, \texttt{idMetAt} \(\}\) \\ \hline
\texttt{Sentencia} & \(\{\) \texttt{;}, \texttt{while}, \texttt{for}, \texttt{if}, \texttt{ret}, \texttt{(}, \texttt{'\{'}, \texttt{id}, \texttt{self}, \texttt{'\}'}, \texttt{else} \(\}\) \\ \hline
\texttt{Asignación} & \(\{\) \texttt{;} \(\}\) \\ \hline
\texttt{Sentencia-Simple} & \(\{\) \texttt{;} \(\}\) \\ \hline
\texttt{Expresión} & \(\{\) \texttt{)}, \texttt{;}, \texttt{]}, \texttt{','} \(\}\) \\ \hline
\texttt{Bloque} & \(\{\) \texttt{;}, \texttt{while}, \texttt{for}, \texttt{if}, \texttt{ret}, \texttt{(}, \texttt{'\{'}, \texttt{id}, \texttt{self}, \texttt{'\}'}, \texttt{else} \(\}\) \\ \hline
\texttt{SentenciaF1} & \(\{\) \texttt{;}, \texttt{while}, \texttt{for}, \texttt{if}, \texttt{ret}, \texttt{(}, \texttt{'\{'}, \texttt{id}, \texttt{self}, \texttt{'\}'}, \texttt{else} \(\}\) \\ \hline
\texttt{SentenciaF2} & \(\{\) \texttt{;}, \texttt{while}, \texttt{for}, \texttt{if}, \texttt{ret}, \texttt{(}, \texttt{'\{'}, \texttt{id}, \texttt{self}, \texttt{'\}'}, \texttt{else} \(\}\) \\ \hline
\texttt{AccesoVar-Simple} & \(\{\) \texttt{=} \(\}\) \\ \hline
\texttt{AccesoSelf-Simple} & \(\{\) \texttt{=} \(\}\) \\ \hline
\texttt{AccesoVar-SimpleF} & \(\{\) \texttt{=} \(\}\) \\ \hline
\texttt{Encadenado-SimpleK} & \(\{\) \texttt{=} \(\}\) \\ \hline
\texttt{Encadenado-Simple} & \(\{\) \texttt{.}, \texttt{=} \(\}\) \\ \hline
\texttt{ExpOr} & \(\{\) \texttt{)}, \texttt{;}, \texttt{]}, \texttt{','} \(\}\) \\ \hline
\texttt{ExpAnd} & \(\{\) \texttt{||}, \texttt{)}, \texttt{;}, \texttt{]}, \texttt{','} \(\}\) \\ \hline
\texttt{ExpOrR} & \(\{\) \texttt{)}, \texttt{;}, \texttt{]}, \texttt{','} \(\}\) \\ \hline
\texttt{ExpIgual} & \(\{\) \texttt{\&\&}, \texttt{||}, \texttt{)}, \texttt{;}, \texttt{]}, \texttt{','} \(\}\) \\ \hline
\texttt{ExpAndR} & \(\{\) \texttt{||}, \texttt{)}, \texttt{;}, \texttt{]}, \texttt{','} \(\}\) \\ \hline
\texttt{ExpCompuesta} & \(\{\) \texttt{==}, \texttt{!=}, \texttt{\&\&}, \texttt{||}, \texttt{)}, \texttt{;}, \texttt{]}, \texttt{','} \(\}\) \\ \hline
\texttt{ExpIgualR} & \(\{\) \texttt{\&\&}, \texttt{||}, \texttt{)}, \texttt{;}, \texttt{]}, \texttt{','} \(\}\) \\ \hline
\texttt{OpIgual} & \(\{\) \texttt{+}, \texttt{-}, \texttt{!}, \texttt{++}, \texttt{-{}-}, \texttt{nil}, \texttt{true}, \texttt{false}, \texttt{intLiteral}, \texttt{StrLiteral}, \texttt{(}, \texttt{self}, \texttt{idclass}, \texttt{new}, \texttt{id} \(\}\) \\ \hline
\texttt{ExpAd} & \(\{\) \texttt{<}, \texttt{>}, \texttt{<=}, \texttt{>=}, \texttt{==}, \texttt{!=}, \texttt{\&\&}, \texttt{||}, \texttt{)}, \texttt{;}, \texttt{]}, \texttt{','} \(\}\) \\ \hline
\texttt{ExpCompuestaF} & \(\{\) \texttt{==}, \texttt{!=}, \texttt{\&\&}, \texttt{||}, \texttt{)}, \texttt{;}, \texttt{]}, \texttt{','} \(\}\) \\ \hline
\texttt{OpCompuesto} & \(\{\) \texttt{+}, \texttt{-}, \texttt{!}, \texttt{++}, \texttt{-{}-}, \texttt{nil}, \texttt{true}, \texttt{false}, \texttt{intLiteral}, \texttt{StrLiteral}, \texttt{(}, \texttt{self}, \texttt{idclass}, \texttt{new}, \texttt{id} \(\}\) \\ \hline
\texttt{ExpMul} & \(\{\) \texttt{+}, \texttt{-}, \texttt{<}, \texttt{>}, \texttt{<=}, \texttt{>=}, \texttt{==}, \texttt{!=}, \texttt{\&\&}, \texttt{||}, \texttt{)}, \texttt{;}, \texttt{]}, \texttt{','} \(\}\) \\ \hline
\texttt{ExpAdR} & \(\{\) \texttt{<}, \texttt{>}, \texttt{<=}, \texttt{>=}, \texttt{==}, \texttt{!=}, \texttt{\&\&}, \texttt{||}, \texttt{)}, \texttt{;}, \texttt{]}, \texttt{','} \(\}\) \\ \hline
\texttt{OpAd} & \(\{\) \texttt{+}, \texttt{-}, \texttt{!}, \texttt{++}, \texttt{-{}-}, \texttt{nil}, \texttt{true}, \texttt{false}, \texttt{intLiteral}, \texttt{StrLiteral}, \texttt{(}, \texttt{self}, \texttt{idclass}, \texttt{new}, \texttt{id} \(\}\) \\ \hline
\texttt{ExpUn} & \(\{\) \texttt{*}, \texttt{/}, \texttt{+}, \texttt{-}, \texttt{<}, \texttt{>}, \texttt{<=}, \texttt{>=}, \texttt{==}, \texttt{!=}, \texttt{\&\&}, \texttt{||}, \texttt{)}, \texttt{;}, \texttt{]}, \texttt{','} \(\}\) \\ \hline
\texttt{ExpMulR} & \(\{\) \texttt{+}, \texttt{-}, \texttt{<}, \texttt{>}, \texttt{<=}, \texttt{>=}, \texttt{==}, \texttt{!=}, \texttt{\&\&}, \texttt{||}, \texttt{)}, \texttt{;}, \texttt{]}, \texttt{','} \(\}\) \\ \hline
\texttt{OpMul} & \(\{\) \texttt{+}, \texttt{-}, \texttt{!}, \texttt{++}, \texttt{-{}-}, \texttt{nil}, \texttt{true}, \texttt{false}, \texttt{intLiteral}, \texttt{StrLiteral}, \texttt{(}, \texttt{self}, \texttt{idclass}, \texttt{new}, \texttt{id} \(\}\) \\ \hline
\texttt{OpUnario} & \(\{\) \texttt{+}, \texttt{-}, \texttt{!}, \texttt{++}, \texttt{-{}-}, \texttt{nil}, \texttt{true}, \texttt{false}, \texttt{intLiteral}, \texttt{StrLiteral}, \texttt{(}, \texttt{self}, \texttt{idclass}, \texttt{new}, \texttt{id} \(\}\) \\ \hline
\texttt{Operando} & \(\{\) \texttt{*}, \texttt{/}, \texttt{+}, \texttt{-}, \texttt{<}, \texttt{>}, \texttt{<=}, \texttt{>=}, \texttt{==}, \texttt{!=}, \texttt{\&\&}, \texttt{||}, \texttt{)}, \texttt{;}, \texttt{]}, \texttt{','} \(\}\) \\ \hline
\texttt{Literal} & \(\{\) \texttt{*}, \texttt{/}, \texttt{+}, \texttt{-}, \texttt{<}, \texttt{>}, \texttt{<=}, \texttt{>=}, \texttt{==}, \texttt{!=}, \texttt{\&\&}, \texttt{||}, \texttt{)}, \texttt{;}, \texttt{]}, \texttt{','} \(\}\) \\ \hline
\texttt{Primario} & \(\{\) \texttt{*}, \texttt{/}, \texttt{+}, \texttt{-}, \texttt{<}, \texttt{>}, \texttt{<=}, \texttt{>=}, \texttt{==}, \texttt{!=}, \texttt{\&\&}, \texttt{||}, \texttt{)}, \texttt{;}, \texttt{]}, \texttt{','} \(\}\) \\ \hline
\texttt{ExpresionParentizada} & \(\{\) \texttt{*}, \texttt{/}, \texttt{+}, \texttt{-}, \texttt{<}, \texttt{>}, \texttt{<=}, \texttt{>=}, \texttt{==}, \texttt{!=}, \texttt{\&\&}, \texttt{||}, \texttt{)}, \texttt{;}, \texttt{]}, \texttt{','} \(\}\) \\ \hline
\texttt{AccesoSelf} & \(\{\) \texttt{*}, \texttt{/}, \texttt{+}, \texttt{-}, \texttt{<}, \texttt{>}, \texttt{<=}, \texttt{>=}, \texttt{==}, \texttt{!=}, \texttt{\&\&}, \texttt{||}, \texttt{)}, \texttt{;}, \texttt{]}, \texttt{','} \(\}\) \\ \hline
\texttt{Llamada-Método-Estático} & \(\{\) \texttt{*}, \texttt{/}, \texttt{+}, \texttt{-}, \texttt{<}, \texttt{>}, \texttt{<=}, \texttt{>=}, \texttt{==}, \texttt{!=}, \texttt{\&\&}, \texttt{||}, \texttt{)}, \texttt{;}, \texttt{]}, \texttt{','} \(\}\) \\ \hline
\texttt{Llamada-Conclassor} & \(\{\) \texttt{*}, \texttt{/}, \texttt{+}, \texttt{-}, \texttt{<}, \texttt{>}, \texttt{<=}, \texttt{>=}, \texttt{==}, \texttt{!=}, \texttt{\&\&}, \texttt{||}, \texttt{)}, \texttt{;}, \texttt{]}, \texttt{','} \(\}\) \\ \hline
\texttt{VarOMetodo} & \(\{\) \texttt{*}, \texttt{/}, \texttt{+}, \texttt{-}, \texttt{<}, \texttt{>}, \texttt{<=}, \texttt{>=}, \texttt{==}, \texttt{!=}, \texttt{\&\&}, \texttt{||}, \texttt{)}, \texttt{;}, \texttt{]}, \texttt{','} \(\}\) \\ \hline
\texttt{AccesoVar} & \(\{\) \texttt{*}, \texttt{/}, \texttt{+}, \texttt{-}, \texttt{<}, \texttt{>}, \texttt{<=}, \texttt{>=}, \texttt{==}, \texttt{!=}, \texttt{\&\&}, \texttt{||}, \texttt{)}, \texttt{;}, \texttt{]}, \texttt{','} \(\}\) \\ \hline
\texttt{ExpresionParentizadaF} & \(\{\) \texttt{*}, \texttt{/}, \texttt{+}, \texttt{-}, \texttt{<}, \texttt{>}, \texttt{<=}, \texttt{>=}, \texttt{==}, \texttt{!=}, \texttt{\&\&}, \texttt{||}, \texttt{)}, \texttt{;}, \texttt{]}, \texttt{','} \(\}\) \\ \hline
\texttt{AccesoSelfF} & \(\{\) \texttt{*}, \texttt{/}, \texttt{+}, \texttt{-}, \texttt{<}, \texttt{>}, \texttt{<=}, \texttt{>=}, \texttt{==}, \texttt{!=}, \texttt{\&\&}, \texttt{||}, \texttt{)}, \texttt{;}, \texttt{]}, \texttt{','} \(\}\) \\ \hline
\texttt{AccesoVarF} & \(\{\) \texttt{*}, \texttt{/}, \texttt{+}, \texttt{-}, \texttt{<}, \texttt{>}, \texttt{<=}, \texttt{>=}, \texttt{==}, \texttt{!=}, \texttt{\&\&}, \texttt{||}, \texttt{)}, \texttt{;}, \texttt{]}, \texttt{','} \(\}\) \\ \hline
\texttt{Llamada-MétodoF} & \(\{\) \texttt{*}, \texttt{/}, \texttt{+}, \texttt{-}, \texttt{<}, \texttt{>}, \texttt{<=}, \texttt{>=}, \texttt{==}, \texttt{!=}, \texttt{\&\&}, \texttt{||}, \texttt{)}, \texttt{;}, \texttt{]}, \texttt{','} \(\}\) \\ \hline
\texttt{Llamada-Método} & \(\{\) \texttt{*}, \texttt{/}, \texttt{+}, \texttt{-}, \texttt{<}, \texttt{>}, \texttt{<=}, \texttt{>=}, \texttt{==}, \texttt{!=}, \texttt{\&\&}, \texttt{||}, \texttt{)}, \texttt{;}, \texttt{]}, \texttt{','} \(\}\) \\ \hline
\texttt{Encadenado} & \(\{\) \texttt{*}, \texttt{/}, \texttt{+}, \texttt{-}, \texttt{<}, \texttt{>}, \texttt{<=}, \texttt{>=}, \texttt{==}, \texttt{!=}, \texttt{\&\&}, \texttt{||}, \texttt{)}, \texttt{;}, \texttt{]}, \texttt{','} \(\}\) \\ \hline
\texttt{Argumentos-Actuales} & \(\{\) \texttt{*}, \texttt{/}, \texttt{+}, \texttt{-}, \texttt{<}, \texttt{>}, \texttt{<=}, \texttt{>=}, \texttt{==}, \texttt{!=}, \texttt{\&\&}, \texttt{||}, \texttt{)}, \texttt{;}, \texttt{]}, \texttt{','}, \texttt{.} \(\}\) \\ \hline
\texttt{Llamada-ConclassorF1} & \(\{\) \texttt{*}, \texttt{/}, \texttt{+}, \texttt{-}, \texttt{<}, \texttt{>}, \texttt{<=}, \texttt{>=}, \texttt{==}, \texttt{!=}, \texttt{\&\&}, \texttt{||}, \texttt{)}, \texttt{;}, \texttt{]}, \texttt{','} \(\}\) \\ \hline
\texttt{Llamada-ConclassorF2} & \(\{\) \texttt{*}, \texttt{/}, \texttt{+}, \texttt{-}, \texttt{<}, \texttt{>}, \texttt{<=}, \texttt{>=}, \texttt{==}, \texttt{!=}, \texttt{\&\&}, \texttt{||}, \texttt{)}, \texttt{;}, \texttt{]}, \texttt{','} \(\}\) \\ \hline
\texttt{Argumentos-ActualesF} & \(\{\) \texttt{*}, \texttt{/}, \texttt{+}, \texttt{-}, \texttt{<}, \texttt{>}, \texttt{<=}, \texttt{>=}, \texttt{==}, \texttt{!=}, \texttt{\&\&}, \texttt{||}, \texttt{)}, \texttt{;}, \texttt{]}, \texttt{','}, \texttt{.} \(\}\) \\ \hline
\texttt{Lista-Expresiones} & \(\{\) \texttt{)} \(\}\) \\ \hline
\texttt{Lista-ExpresionesF} & \(\{\) \texttt{)} \(\}\) \\ \hline
\texttt{EncadenadoF} & \(\{\) \texttt{*}, \texttt{/}, \texttt{+}, \texttt{-}, \texttt{<}, \texttt{>}, \texttt{<=}, \texttt{>=}, \texttt{==}, \texttt{!=}, \texttt{\&\&}, \texttt{||}, \texttt{)}, \texttt{;}, \texttt{]}, \texttt{','} \(\}\) \\ \hline
\end{longtable}
\endgroup


\section{Análisis de la condición LL(1)}
Como se mencionó en la sección~\ref{sec:eliminacion-ambiguedad},
la gramática final no es LL(1) debido a la ambigüedad de las
sentencias \texttt{if} y \texttt{else}.

El analizador sintáctico predictivo descendente recursivo resuelve
este conflicto mediante la siguiente regla: cuando el
token es \texttt{else}, se lo consume y se lo
asocia al \texttt{if} más cercano que todavía no tiene un
\texttt{else}. Esta decisión permite analizar los condicionales
anidados de forma determinista, aunque la gramática siga sin
ser LL(1).

\section{Diseño e implementación}
\begin{description}[style=nextline]
  \item[\lexcode{SyntacticAnalyzer}]
  Conserva el lector léxico y el token actual. Cada no terminal se
  corresponde con un método que reconoce su producción y consume tokens.
  \item[\lexcode{Lexer} y \lexcode{Token}]
  Proporcionan la secuencia de entrada y las posiciones necesarias para
  informar dónde aparece un símbolo inesperado.
  \item[\lexcode{Etapa2}]
  Recibe el archivo fuente, construye los analizadores y ejecuta
  \lexcode{program()}. Las instrucciones comunes de compilación y ejecución
  se encuentran en la sección~\ref{sec:compilacion-ejecucion}.
\end{description}

El método \lexcode{matchTokenType()} compara el tipo del token actual con el
esperado. Si coinciden, solicita el siguiente token; en caso contrario,
lanza una excepción. Los métodos \lexcode{checkType()} y
\lexcode{checkLexeme()} permiten seleccionar alternativas sin consumir
anticipadamente la entrada. La comprobación final de EOF impide aceptar
tokens adicionales después del bloque \texttt{start}.

El lexema \texttt{start} se recibe como identificador de variable o método.
El método \lexcode{start()} comprueba su texto para distinguirlo de
otros identificadores en esa posición del programa.

\section{Errores sintácticos}
\label{sec:errores-sintacticos}
\begingroup
\setlength{\emergencystretch}{3em}
El analizador interrumpe el reconocimiento al detectar el primer error.
La ubicación informada corresponde al token actual recibido del analizador
léxico: indica dónde se detectó la incompatibilidad con la gramática, que
puede ser posterior al símbolo omitido. No se realiza recuperación de errores.

\subsection{Excepciones utilizadas}
Las clases pertenecen al paquete
\path{uncuyo.compiladores.exceptions.syntactic}.
\begin{description}[style=nextline]
  \item[\lexcode{SyntacticException}]
  Hereda de \lexcode{Exception}. Su constructor recibe la posición
  (\lexcode{rowIndex}, \lexcode{columnIndex}) y la descripción
  (\lexcode{description}) para construir el mensaje.
  Es la excepción general utilizada cuando falta un token o
  ninguna alternativa de una producción admite la entrada actual.

  \item[\lexcode{MethodOutsideImplException}]
  Hereda de \lexcode{SyntacticException} y recibe un \lexcode{Token}.
  Se lanza desde \lexcode{start()} cuando el lexema actual no es
  \texttt{start}. Su descripción es
  \texttt{El metodo '<lexema>' debe estar contenido dentro de un impl}.
  Por ejemplo, \lexcode{invalidMethod.s} informa este error para
  \texttt{m1}, en la línea 5, columna 1. La comprobación se basa en
  el lexema encontrado al entrar en \lexcode{start()}, por lo que esta
  excepción no representa todos los errores de métodos o constructores
  fuera de \texttt{impl}.
\end{description}

Las excepciones \lexcode{LexicalException} y
\lexcode{LexerFileReadException} provienen de la lectura léxica y conservan
sus propios diagnósticos. No se convierten en excepciones sintácticas.

\subsection{Cómo se detecta el error}
El método \lexcode{matchTokenType()} verifica un token concreto. Si el
tipo actual no coincide con el esperado, lanza una
\lexcode{SyntacticException} con el patrón
\texttt{Se esperaba <esperado>, pero se encontro '<encontrado>'}.
El valor esperado se obtiene mediante \lexcode{TokenDefinitions.getLexeme()}:
si existe un lexema asociado al tipo, se muestra ese símbolo; en caso
contrario, se muestra el nombre del tipo, como
\lexcode{IDENTIFIER_CLASS}, \lexcode{IDENTIFIER_VAR_METHOD} o \lexcode{EOF}.

Los métodos que seleccionan alternativas de una producción construyen
descripciones con varios símbolos posibles. Por ejemplo,
\lexcode{classF()} admite una llave de apertura o dos puntos;
\lexcode{miembroK()} admite \texttt{st}, \texttt{fn}, un punto o una llave
de cierre; y \lexcode{expUn()} comprueba los posibles inicios de una
expresión. Cuando el token no pertenece a las alternativas admitidas,
la excepción enumera los símbolos esperados y muestra el lexema encontrado.
La redacción de esas listas es la que está escrita en cada método.

Al finalizar \lexcode{program()}, se exige \lexcode{EOF}. Por eso,
\lexcode{trailingTokensAfterStart.s} produce
\texttt{Se esperaba EOF, pero se encontro 'class'} en la línea 2,
columna 1. Si falta \texttt{start} y se alcanza el fin de archivo,
\lexcode{listaDefiniciones()} informa
\texttt{Se esperaba 'class', 'impl' o 'start', pero se encontro 'EOF'}.

\subsection{Formato del mensaje}
El método \lexcode{getMessage()} de \lexcode{SyntacticException} devuelve
un encabezado y una fila con la ubicación y la descripción. Por ejemplo,
el caso \lexcode{lowerCaseClass.s} genera:
\begin{lstlisting}[style=bnftinys]
ERROR SINTACTICO
| NÚMERO DE LÍNEA (NÚMERO DE COLUMNA) | DESCRIPCIÓN: |
| LÍNEA 1 (COLUMNA 7) | Se esperaba IDENTIFIER_CLASS, pero se encontro 'miClase' |
\end{lstlisting}
La descripción conserva la redacción de la implementación, incluidas
las diferencias de puntuación y los nombres internos de tipos de token.

En \lexcode{Etapa2}, \lexcode{printStackTrace()} muestra las excepciones
sintácticas por la salida de error. Esta salida incluye el nombre
completo de la clase de excepción, el mensaje anterior y la traza de
llamadas de Java. Los errores léxicos que se capturan durante
\lexcode{program()} se imprimen con el prefijo
\texttt{ERROR: LEXICO}, seguido del mensaje léxico. Si el análisis
termina correctamente, se informa:
\begin{verbatim}
CORRECTO: ANALISIS SINTACTICO
\end{verbatim}
La implementación utiliza \texttt{ERROR SINTACTICO} como encabezado.
Para la entrega se exige \texttt{ERROR: SINTACTICO}.
La presentación actual también incluye la traza de Java.

\endgroup

\clearpage
\section{Escenarios y casos de prueba}
\begingroup
\setlength{\emergencystretch}{3em}
Los archivos se encuentran en \lexcode{tc/etapa2}, dentro de la raíz del
proyecto, y se distribuyen en dos directorios:
\begin{itemize}[leftmargin=*]
  \item Casos válidos: \lexcode{tc/etapa2/success}, con 29 archivos de prueba.
  \item Casos inválidos: \lexcode{tc/etapa2/error}, con 30 archivos de prueba.
\end{itemize}
Los escenarios cubren la estructura del programa, clases e implementaciones,
herencia, atributos, métodos, constructores, arreglos, expresiones y
sentencias de control. Esta organización utiliza el directorio
\lexcode{tc} solicitado por el enunciado y separa los casos por etapa.

\subsection{Técnica y alcance de la verificación}
La clase \lexcode{SyntacticTest} agrupa las pruebas de esta etapa.
Los casos válidos usan \lexcode{assertDoesNotThrow} para comprobar que
\lexcode{program()} alcance EOF.

Los casos inválidos usan \lexcode{assertThrows} para exigir una
\lexcode{SyntacticException} o una subclase, como
\lexcode{MethodOutsideImplException}. No se comparan el texto del
mensaje ni su ubicación.
Desde la raíz del proyecto se ejecutan mediante:
\begin{verbatim}
mvn -Dtest=SyntacticTest test
\end{verbatim}

\subsection{Casos que deben aceptarse}
En cada caso se espera alcanzar EOF sin excepciones.
\begingroup
\small
\renewcommand{\arraystretch}{1.2}
\setlength{\tabcolsep}{5pt}
\setlength{\LTleft}{0pt}
\setlength{\LTright}{0pt plus 1fill}
\begin{longtable}{|
>{\raggedright\arraybackslash}p{0.40\textwidth}|
>{\raggedright\arraybackslash}p{0.28\textwidth}|
>{\raggedright\arraybackslash}p{\dimexpr0.32\textwidth-6\tabcolsep-4\arrayrulewidth\relax}|}
\caption{Casos válidos del analizador sintáctico}\label{tab:sintactico-validos}\\
\hline
\textbf{Archivo} & \textbf{Propósito} & \textbf{Resultado esperado} \\ \hline
\endfirsthead
\hline
\textbf{Archivo} & \textbf{Propósito} & \textbf{Resultado esperado} \\ \hline
\endhead
\texttt{array\allowbreak .s} & Reconocer arreglos, índices y un recorrido con for. & Aceptar creación, asignaciones y recorrido. \\ \hline
\texttt{attributeAccess\allowbreak .s} & Reconocer asignaciones a atributos de un objeto. & Aceptar el acceso a.x. \\ \hline
\texttt{chainedMethod\allowbreak .s} & Reconocer llamadas y atributos encadenados. & Aceptar a.crearB(7).val. \\ \hline
\texttt{class\&start\allowbreak .s} & Reconocer una clase seguida del bloque principal. & Aceptar la sintaxis sin exigir impl. \\ \hline
\texttt{classA\&B\allowbreak .s} & Reconocer varias clases, bloques impl y accesos encadenados. & Aceptar la estructura completa. \\ \hline
\texttt{classA\allowbreak .s} & Reconocer un método que retorna otra clase. & Aceptar la sintaxis sin resolver el tipo B. \\ \hline
\texttt{classesABC\allowbreak .s} & Reconocer varias clases y una declaración de herencia. & Aceptar las definiciones antes de start. \\ \hline
\texttt{classPersona\&impl\allowbreak .s} & Reconocer atributos, un método y su llamada. & Aceptar el método sin tipo de retorno explícito. \\ \hline
\texttt{classPersona\allowbreak .s} & Reconocer una clase con dos atributos Str. & Aceptar la clase y el bloque principal vacío. \\ \hline
\texttt{ejemploManual1\allowbreak .s} & Reconocer el ejemplo de arreglos del manual. & Aceptar índices, for y llamadas a IO. \\ \hline
\texttt{ejemploManual2\allowbreak .s} & Reconocer el ejemplo de Fibonacci del manual. & Aceptar métodos, while y condicionales anidados. \\ \hline
\texttt{emptyStart\allowbreak .s} & Reconocer sentencias vacías dentro de start. & Aceptar los puntos y coma aislados. \\ \hline
\texttt{herencia\allowbreak .s} & Reconocer la herencia indicada mediante dos puntos. & Aceptar la clase base y la derivada. \\ \hline
\texttt{if\allowbreak .s} & Reconocer condicionales anidados sin else. & Aceptar las sentencias sin llaves. \\ \hline
\texttt{ifElse\allowbreak .s} & Reconocer un condicional con alternativa. & Aceptar las ramas if y else. \\ \hline
\texttt{ioOperations\allowbreak .s} & Reconocer una llamada estática a IO. & Aceptar la llamada como sentencia entre paréntesis. \\ \hline
\texttt{methodLocalVariables\allowbreak .s} & Reconocer variables locales al inicio de un método. & Aceptar declaraciones, asignaciones y retorno. \\ \hline
\texttt{methodWithParams\allowbreak .s} & Reconocer parámetros formales y argumentos actuales. & Aceptar las listas separadas por comas. \\ \hline
\texttt{methodWithReturn\allowbreak .s} & Reconocer un método Int y su sentencia ret. & Aceptar el retorno de un literal entero. \\ \hline
\texttt{motor\allowbreak .s} & Reconocer un atributo Array y un método void con for. & Aceptar el recorrido y el if con else. \\ \hline
\texttt{multipleConstructors\allowbreak .s} & Reconocer dos constructores en un mismo impl. & Aceptar su sintaxis; la duplicación es semántica. \\ \hline
\texttt{multipleImpl\allowbreak .s} & Reconocer varios bloques impl de una clase. & Aceptar los métodos distribuidos entre bloques. \\ \hline
\texttt{nestedControlAndEmptySentences\allowbreak .s} & Combinar while, if, else y sentencias vacías. & Aceptar el anidamiento y el incremento. \\ \hline
\texttt{operaciones\allowbreak .s} & Reconocer métodos de instancia y estáticos. & Aceptar st fn y las llamadas anidadas. \\ \hline
\texttt{paramConstructor\allowbreak .s} & Reconocer un constructor con dos parámetros. & Aceptar las asignaciones mediante self. \\ \hline
\texttt{retMethod\allowbreak .s} & Reconocer un método con parámetros y retorno. & Aceptar la suma en la sentencia ret. \\ \hline
\texttt{start\allowbreak .s} & Reconocer el programa mínimo. & Alcanzar EOF después de start vacío. \\ \hline
\texttt{unaryOperators\allowbreak .s} & Reconocer incremento y decremento prefijos. & Aceptar los operadores en asignaciones. \\ \hline
\texttt{variables\allowbreak .s} & Reconocer declaración y asignaciones de una variable. & Aceptar asignaciones con operadores unarios. \\ \hline
\end{longtable}
\endgroup


\subsection{Casos que deben rechazarse}
En cada caso se espera una \lexcode{SyntacticException} por el motivo
indicado. En \lexcode{invalidMethod.s}, la excepción concreta es
\lexcode{MethodOutsideImplException}.
Se describe el primer error detectado en cada archivo, aunque su nombre
sugiera otro problema posterior.
\begingroup
\small
\renewcommand{\arraystretch}{1.2}
\setlength{\tabcolsep}{5pt}
\setlength{\LTleft}{0pt}
\setlength{\LTright}{0pt plus 1fill}
\begin{longtable}{|
>{\raggedright\arraybackslash}p{0.40\textwidth}|
>{\raggedright\arraybackslash}p{0.28\textwidth}|
>{\raggedright\arraybackslash}p{\dimexpr0.32\textwidth-6\tabcolsep-4\arrayrulewidth\relax}|}
\caption{Casos inválidos del analizador sintáctico}\label{tab:sintactico-invalidos}\\
\hline
\textbf{Archivo} & \textbf{Propósito} & \textbf{Resultado esperado} \\ \hline
\endfirsthead
\hline
\textbf{Archivo} & \textbf{Propósito} & \textbf{Resultado esperado} \\ \hline
\endhead
\texttt{classAfterStart\allowbreak .s} & Rechazar un acceso a miembro sobre un literal entero. & Error en el punto de 10.out\_str, antes de la clase final. \\ \hline
\texttt{constructorOutsideImpl\allowbreak .s} & Rechazar un constructor fuera de impl. & Error ante el punto en el nivel superior. \\ \hline
\texttt{declarationInsideBlock\allowbreak .s} & Rechazar una declaración dentro de un bloque if. & Error ante Int en el bloque de sentencias. \\ \hline
\texttt{herenciaMultiple\allowbreak .s} & Rechazar una lista de clases base. & Error: se esperaba una llave en lugar de la coma. \\ \hline
\texttt{inccorectInheritanceDeclaration\allowbreak .s} & Rechazar extends como notación de herencia. & Error: se esperaba una llave o dos puntos. \\ \hline
\texttt{incorrectArrayAccess\allowbreak .s} & Rechazar paréntesis usados como índice de un arreglo. & Error ante el paréntesis de a(0). \\ \hline
\texttt{incorrectForCondition\allowbreak .s} & Rechazar dos puntos en lugar de in en un for. & Error: se esperaba in. \\ \hline
\texttt{incorrectInitialization\allowbreak .s} & Rechazar la declaración Int[] datos. & Error: se esperaba un identificador tras Int. \\ \hline
\texttt{incorrectWhileCondition\allowbreak .s} & Rechazar una condición while sin paréntesis. & Error: se esperaba el paréntesis de apertura. \\ \hline
\texttt{invalidConstructor\allowbreak .s} & Rechazar un constructor escrito con el nombre de la clase. & Error ante MiClase como miembro de impl. \\ \hline
\texttt{invalidMethod\allowbreak .s} & Rechazar un método fuera de impl. & Error específico de método fuera de impl para m1. \\ \hline
\texttt{lowerCaseClass\allowbreak .s} & Rechazar un nombre de clase que comienza con minúscula. & Error: se esperaba un identificador de clase. \\ \hline
\texttt{methodWithoutBody\allowbreak .s} & Rechazar un método terminado en punto y coma. & Error: se esperaba la llave de apertura del cuerpo. \\ \hline
\texttt{methodWithoutFn\allowbreak .s} & Rechazar un método sin la palabra fn. & Error ante Int en la lista de miembros. \\ \hline
\texttt{multipleComparisons\allowbreak .s} & Rechazar comparaciones relacionales encadenadas. & Error ante el segundo operador menor que. \\ \hline
\texttt{multipleInheritance\allowbreak .s} & Rechazar dos clases base separadas por coma. & Error: se esperaba una llave en lugar de la coma. \\ \hline
\texttt{newDePrimitivo\allowbreak .s} & Rechazar new Int() como creación de un objeto. & Error: se esperaba un corchete para crear un arreglo. \\ \hline
\texttt{noStartMethod\allowbreak .s} & Detectar un programa sin bloque start. & Error al encontrar EOF entre las definiciones. \\ \hline
\texttt{noTypeAssignation\allowbreak .s} & Rechazar una creación de arreglo sin tipo. & Error ante el corchete posterior a new. \\ \hline
\texttt{noTypeParameter\allowbreak .s} & Rechazar un parámetro formal sin tipo. & Error ante a en la lista de parámetros. \\ \hline
\texttt{pubMethod\allowbreak .s} & Rechazar pub como modificador de un método. & Error ante pub en la lista de miembros. \\ \hline
\texttt{redefinirPrimitivo\allowbreak .s} & Rechazar un tipo primitivo como nombre de clase. & Error: Int no es un identificador de clase. \\ \hline
\texttt{selfAsVariableName\allowbreak .s} & Rechazar self como nombre de un atributo. & Error: se esperaba un identificador de variable. \\ \hline
\texttt{trailingTokensAfterStart\allowbreak .s} & Rechazar definiciones posteriores al bloque principal. & Error: se esperaba EOF y se encontró class. \\ \hline
\texttt{unclosedClass\allowbreak .s} & Detectar una clase sin llave de cierre. & Error ante impl dentro de la lista de atributos. \\ \hline
\texttt{unclosedIf\allowbreak .s} & Rechazar la inicialización dentro de una declaración local. & Error en Int a = 1; antes de llegar al if incompleto. \\ \hline
\texttt{unclosedLines\allowbreak .s} & Detectar la falta de punto y coma tras una declaración. & Error al encontrar x en lugar del punto y coma. \\ \hline
\texttt{UpperCaseVariable\allowbreak .s} & Rechazar una variable que comienza con mayúscula. & Error: MiVar se reconoce como nombre de clase. \\ \hline
\texttt{variableDeclarationOnSentence\allowbreak .s} & Rechazar una declaración local dentro de un if. & Error ante Int en el bloque de sentencias. \\ \hline
\texttt{wrongMethodCall\allowbreak .s} & Rechazar una llamada usada sin paréntesis de sentencia. & Error ante el paréntesis de io.out\_str(...). \\ \hline
\end{longtable}
\endgroup

\endgroup
